% =====================================================================
% sec/3_method.tex  --  Methodology
% =====================================================================

\section{Methodology}
\label{sec:method}

\subsection{Overview}
\label{subsec:overview}

\cref{fig:arch} illustrates the overall ProbVLP architecture. Given an input
endoscopic image $\bm{x} \in \mathbb{R}^{3 \times H \times W}$ and a text prompt
$\bm{t}$ (e.g., ``\textit{a colonoscopy image of a polyp}''), ProbVLP produces
a segmentation logit map $\hat{\bm{y}} \in \mathbb{R}^{H \times W}$ and a
calibrated uncertainty map $\bm{u} \in \mathbb{R}^{H \times W}$.

\subsection{EndoViT Visual Backbone}
\label{subsec:endovit}

We replace the CLIP ViT-B/16 visual tower with \textbf{EndoViT}, a ViT-B/16
pretrained via masked autoencoding (MAE) on Endo700k -- a large-scale collection
of endoscopic images. EndoViT produces a sequence of patch tokens
$\bm{F}_v \in \mathbb{R}^{(1+N_p) \times d}$, where $N_p = (H/p)^2$ for patch
size $p{=}16$ and $d{=}768$.

\textbf{Domain renormalization.} Because the incoming image is normalized to CLIP
statistics $(\mu_\text{CLIP}, \sigma_\text{CLIP})$ by the upstream dataloader, we
undo this normalization and renormalize to EndoViT's pretraining statistics
$(\mu_\text{EndoViT}, \sigma_\text{EndoViT})$ at the patch-embedding stage:
\begin{equation}
    \hat{\bm{x}} = \frac{\bm{x} \odot \sigma_\text{CLIP} + \mu_\text{CLIP}
                   - \mu_\text{EndoViT}}{\sigma_\text{EndoViT}}.
\end{equation}

\subsection{LoRA-Adapted PVL Adapters}
\label{subsec:lora}

The PVL adapters~\cite{} use bidirectional cross-attention to fuse image and
text tokens. Let $\bm{W}_0 \in \mathbb{R}^{d_\text{out} \times d_\text{in}}$
denote a frozen projection matrix (query, key, value, or output). We augment it
with a LoRA delta:
\begin{equation}
    \bm{W} = \bm{W}_0 + \frac{\alpha}{r}\,\bm{B}\bm{A},
\end{equation}
where $\bm{A} \in \mathbb{R}^{r \times d_\text{in}}$,
$\bm{B} \in \mathbb{R}^{d_\text{out} \times r}$, $r{=}8$, $\alpha{=}16$.
$\bm{A}$ is initialized with Kaiming uniform, $\bm{B}$ with zeros, so the delta
is zero at initialization and training starts from the pretrained MedCLIPSeg
weights. We apply LoRA to $\{$\texttt{q\_proj, k\_proj, v\_proj, out\_proj}$\}$
in both cross-attention directions (\emph{image-to-text} and \emph{text-to-image}).

\subsection{Boundary-Feedback Gating}
\label{subsec:boundary}

A lightweight boundary head $\mathcal{B}$ predicts boundary logits
$\bm{b}^{(t)} \in \mathbb{R}^{H \times W}$ from the upscaled feature map at
training step $t$. At step $t{+}1$, a downsampled version of $\sigma(\bm{b}^{(t)})$
biases the image token stream entering each PVL adapter:
\begin{equation}
    \tilde{\bm{f}}_i = \bm{f}_i + \tanh(g_i)\,
        \mathcal{P}(\sigma(\bm{b}^{(t)})),
\end{equation}
where $g_i$ is a per-adapter learnable scalar gate (initialized to 0) and
$\mathcal{P}$ denotes adaptive average pooling to the patch grid. $g_i$ starts
at 0 (identical to no gating), allowing the network to learn the appropriate
feedback magnitude from data.

\subsection{Segmentation Head}
\label{subsec:seghead}

The upscaled feature map $\bm{F} \in \mathbb{R}^{C \times H' \times W'}$ is
processed by a Squeeze-and-Excitation (SE) block~\cite{hu2018senet} for
channel recalibration:
\begin{equation}
    \bm{F}_\text{SE} = \bm{F} \odot \sigma\!\left(
        \bm{W}_2 \,\text{GELU}\!\left(\bm{W}_1\,\GAP(\bm{F})\right)\right).
\end{equation}
Segmentation logits are obtained via dot-product similarity with the projected
text embedding:
\begin{equation}
    \hat{\bm{y}} = \text{Up}\!\left(
        \bm{W}_\text{mask}(\hat{\bm{t}})\cdot\bm{F}_\text{SE}\right),
\end{equation}
where $\hat{\bm{t}} = \bm{t}/\norm{\bm{t}}$ is the $\ell_2$-normalized text
feature.

\subsection{Training Objective}
\label{subsec:losses}

We minimize a composite loss:
\begin{equation}
    \mathcal{L} = \mathcal{L}_\text{Dice} + \mathcal{L}_\text{BCE}
                + \lambda_T\,\mathcal{L}_\text{Tversky}
                + \lambda_B\,\mathcal{L}_\text{Boundary},
\end{equation}
where $\lambda_T{=}0.3$, $\lambda_B{=}0.3$. The Tversky loss uses $\alpha{=}0.3$,
$\beta{=}0.7$ to penalize false negatives more heavily. $\mathcal{L}_\text{Boundary}$
is binary cross-entropy between boundary logits and soft boundary labels obtained
via morphological dilation/erosion of the ground-truth mask.

\subsection{Inference with Inverse-Variance TTA}
\label{subsec:tta}

At test time we generate $K{=}6$ augmented views (original $+$ horizontal flip,
each at three scales $\{0.85, 1.0, 1.15\}$), each processed with $M{=}20$
MC-dropout samples. The fused prediction and uncertainty are:
\begin{align}
    \hat{\bm{p}} &= \frac{\sum_k \bm{m}_k / \bm{v}_k}{\sum_k 1 / \bm{v}_k}, \\
    \bm{u}       &= \left(\sum_k 1 / \bm{v}_k\right)^{-1},
\end{align}
where $\bm{m}_k$ and $\bm{v}_k$ are the per-view mean and variance from the
MC-dropout samples. Pixels with $\hat{\bm{p}} > 0.5$ are classified as polyp.
